\documentclass[a4paper,11pt,xelatex,ja=standard]{bxjsarticle}

\geometry{top=20mm,bottom=25mm,left=20mm,right=20mm}
\usepackage{amsmath}
\usepackage{booktabs}
\usepackage{longtable}

\title{今週の進捗報告\\[0.5em]\large EMS最適化手法の比較研究における\\論文執筆に向けた準備と実験結果}
\author{神戸大学大学院システム情報学研究科 修士1年\\山崎博之}
\date{2026年10月1日}

\begin{document}
\maketitle

\section{背景と目的}
本研究では，太陽光発電（PV）と蓄電池を備えたエネルギーマネジメントシステム（EMS）について，電力料金を最小化する充放電計画を対象とする．論文執筆に向け，線形計画法（LP）と動的計画法（DP）を同一条件で比較し，解の費用差，計算時間，メモリ使用量を評価する準備を進めた．本報告では，9月29日までに行ったモデルの整理，実装の検証，比較実験，研究方針の検討をまとめる．

\section{論文執筆に向けた準備}
まず，需要・PV・蓄電池1台・系統買電からなる共通モデルを定めた．対象期間内の最大買電電力に基づく基本料金相当額と，買電量に基づく電力量料金の合計を最小化する．2025年の実データを使用し，電力量単価21.51\,円/kWh，蓄電池容量860\,kWh，最大充放電出力400\,kW，初期・終端SOC50\%などの条件を揃えた．基本料金は対象期間の長さに応じて按分した．

次に，LP，充放電の排他制約を明示した混合整数線形計画法（MILP），およびSOCを離散化したDPを実装した．人工データの3ケースでは，DPと同じSOC格子に制限したMILPの最大買電電力，電力量料金，基本料金，合計料金がすべて一致した．また，計算結果を独立した評価処理で再計算し，電力収支，SOC遷移，充放電制約，料金内訳を確認できるようにした．

その上で，計画期間，時間刻み，季節，蓄電池出力を変えた実験を行った．各条件を4回実行し，費用差と資源使用量は成功試行の中央値で比較した．1回当たりの時間上限は300秒である．本実験で新たに実行した168件は143件成功，25件時間切れ，失敗0件だった．初期実験の再利用分を含む264件のうち，成功した227件はすべて独立評価に合格した．

\section{比較実験の結果}
表\ref{tab:period}--\ref{tab:battery}にDPの結果を示す．費用差は同じ条件の連続SOC・MILP最適費用を基準にした相対差である．時間とメモリは全処理時間およびピークRSSの中央値であり，時間切れの試行から費用差を推定していない．表中の「--」は成功試行がなく，値を評価できないことを示す．

\begingroup
\small
\setlength{\tabcolsep}{5pt}
\begin{longtable}{llrrrr}
\caption{期間・時間刻み・SOC格子によるDPの比較（各4回）}\label{tab:period}\\
\toprule
条件 & 格子点数 & 成功/4 & 費用差 [\%] & 時間 [秒] & メモリ [MiB]\\
\midrule
\endfirsthead
\toprule
条件 & 格子点数 & 成功/4 & 費用差 [\%] & 時間 [秒] & メモリ [MiB]\\
\midrule
\endhead
1日・30分 & 19 & 4 & 9.99 & 0.03 & 0.1\\
 & 37 & 4 & 5.36 & 0.20 & 0.6\\
 & 73 & 4 & 3.08 & 1.7 & 3.3\\
 & 145 & 4 & 1.57 & 12.1 & 20.5\\
\midrule
7日・30分 & 19 & 4 & 5.55 & 0.3 & 1.4\\
 & 37 & 4 & 4.08 & 6.5 & 6.0\\
 & 73 & 4 & 2.10 & 53.6 & 21.5\\
 & 145 & 0 & -- & -- & --\\
\midrule
30日・30分 & 19 & 4 & 4.12 & 9.0 & 7.3\\
 & 37 & 4 & 2.22 & 89.3 & 18.4\\
 & 73 & 0 & -- & -- & --\\
 & 145 & 0 & -- & -- & --\\
\midrule
7日・60分 & 19 & 4 & 2.37 & 14.1 & 1.7\\
 & 37 & 4 & 0.82 & 103.4 & 9.1\\
 & 73 & 3 & 0.40 & 132.5 & 49.0\\
 & 145 & 0 & -- & -- & --\\
\bottomrule
\end{longtable}
\endgroup

1日の30分刻みでは，格子点数を19点から145点へ増やすと費用差は9.99\%から1.57\%へ縮小した一方，処理時間は0.03秒から12.1秒へ増えた．7日・30分刻みでは73点の費用差が2.10\%，処理時間が53.6秒であり，145点は4回とも300秒で時間切れとなった．30日・30分刻みでは73点と145点が全回時間切れとなった．7日・60分刻みの費用差は集約後の60分問題に対する値であり，30分問題の誤差と直接比較しない．

\begingroup
\small
\setlength{\tabcolsep}{5pt}
\begin{longtable}{llrrrr}
\caption{季節別1週間におけるDPの比較（各4回，30分刻み）}\label{tab:season}\\
\toprule
季節 & 格子点数 & 成功/4 & 費用差 [\%] & 時間 [秒] & メモリ [MiB]\\
\midrule
\endfirsthead
\toprule
季節 & 格子点数 & 成功/4 & 費用差 [\%] & 時間 [秒] & メモリ [MiB]\\
\midrule
\endhead
冬 & 19 & 4 & 82.95 & 0.2 & 0.4\\
 & 37 & 4 & 14.96 & 0.8 & 0.5\\
 & 73 & 4 & 7.03 & 11.3 & 5.2\\
 & 145 & 4 & 3.69 & 131.6 & 37.6\\
\midrule
春 & 19 & 4 & 12.11 & 0.2 & 0.2\\
 & 37 & 4 & 6.50 & 2.0 & 2.2\\
 & 73 & 4 & 3.18 & 22.6 & 12.9\\
 & 145 & 4 & 1.28 & 231.5 & 68.0\\
\midrule
夏 & 19 & 4 & 5.55 & 0.3 & 1.4\\
 & 37 & 4 & 4.08 & 6.5 & 6.0\\
 & 73 & 4 & 2.10 & 53.6 & 21.5\\
 & 145 & 0 & -- & -- & --\\
\midrule
秋 & 19 & 4 & 8.71 & 0.5 & 1.3\\
 & 37 & 4 & 3.65 & 6.9 & 5.3\\
 & 73 & 4 & 2.12 & 64.8 & 23.3\\
 & 145 & 0 & -- & -- & --\\
\bottomrule
\end{longtable}
\endgroup

同じ73点格子でも，処理時間は冬11.3秒，春22.6秒，夏53.6秒，秋64.8秒と変化した．145点格子は冬と春では4回とも成功したが，夏と秋では4回とも時間切れとなった．ただし，各季節から選んだ1週間ずつの結果であり，季節全体の代表値とまではいえない．

\begingroup
\small
\setlength{\tabcolsep}{5pt}
\begin{longtable}{llrrrr}
\caption{蓄電池出力条件別1週間におけるDPの比較（各4回，30分刻み）}\label{tab:battery}\\
\toprule
出力条件 & 格子点数 & 成功/4 & 費用差 [\%] & 時間 [秒] & メモリ [MiB]\\
\midrule
\endfirsthead
\toprule
出力条件 & 格子点数 & 成功/4 & 費用差 [\%] & 時間 [秒] & メモリ [MiB]\\
\midrule
\endhead
低（0.25\,h$^{-1}$） & 19 & 4 & 3.66 & 1.0 & 2.6\\
 & 37 & 4 & 2.08 & 8.7 & 7.4\\
 & 73 & 4 & 1.08 & 77.1 & 29.2\\
 & 145 & 0 & -- & -- & --\\
\midrule
基準（0.465\,h$^{-1}$） & 19 & 4 & 5.55 & 0.3 & 1.4\\
 & 37 & 4 & 4.08 & 6.5 & 6.0\\
 & 73 & 4 & 2.10 & 53.6 & 21.5\\
 & 145 & 0 & -- & -- & --\\
\midrule
高（0.75\,h$^{-1}$） & 19 & 4 & 3.66 & 1.8 & 2.7\\
 & 37 & 4 & 2.08 & 16.3 & 8.4\\
 & 73 & 4 & 1.08 & 146.6 & 37.0\\
 & 145 & 0 & -- & -- & --\\
\bottomrule
\end{longtable}
\endgroup

蓄電池容量を860\,kWhに固定して出力／容量比を変えると，73点格子の処理時間は低出力77.1秒，基準53.6秒，高出力146.6秒となった．計算時間は出力比だけでは単調に決まらず，実行時に保持する状態や遷移の数にも左右された．

\section{結果の解釈と論文の方針}
全11条件で成功したLP44件は，対応するMILP基準と費用が一致し，観測された同時充放電は0\,kWだった．費用差の許容幅0.1・0.5・1.0\%，時間上限1・10・60・300秒，メモリ上限64・512・2048\,MiBを組み合わせた396条件では，387条件でLPが選ばれ，残る9条件には全要件を満たす手法がなかった．これらは本研究の固定正価格，売電なし，蓄電池1台という条件に限る結果である．

9月29日の検討では，この実験結果とモデルの条件を踏まえ，論文の中心をLPとDPの比較に置き，MILPは充放電の排他制約を明示した検証用の基準として位置づける方針を整理した．LPで同時充放電が起きなかったという観測結果と，異なる料金・設備条件でも同じことが成り立つという一般的な主張は区別する．

\section{今後の課題}
次は，実験結果を論文原稿に統合し，格子細分化による費用差の縮小と計算負荷の増加，季節・設備条件による違い，LPを選ぶ際の前提を説明する．300秒で時間切れとなった設定の費用と必要時間は未確認である．また，季節ごとに1週間のみを調べたため，季節内の変動や固定価格以外の条件への一般化は今後の検証課題とする．

\end{document}
